%=============================================================
%  SEGMENT 4 — Chapters 13–16
%  Book: وعده یا خدعه؟ — بررسی انتقادی
%  Codename: khomeini-promise-critical-study
%  Requires: Segment 0 preamble compiled; XeLaTeX ×2
%=============================================================

%%%%%%%%%%%%%%%%%%%%%%%%%%%%%%%%%%%%%%%%%%%%%%%%%%%%%%%%%%%%
%%  CHAPTER 13
%%%%%%%%%%%%%%%%%%%%%%%%%%%%%%%%%%%%%%%%%%%%%%%%%%%%%%%%%%%%
\chapter{تحلیل تطبیقی — خمینی در هر بزنگاه چه کرد؟}
\label{ch:comparative}

%------------------------------------------------------------
\begin{infoBox}[چکیدهٔ فصل ۱۳]
این فصل هفت بزنگاه کلیدی (فصل‌های ۶–۱۲) را در یک جدول تطبیقیِ واحد گرد می‌آورد.
برای هر بزنگاه، کنش واقعیِ خمینی با پیش‌بینیِ هر یک از چهار تز سنجیده می‌شود
و میزان سازگاری هر تز با شواهد تجربی ارزیابی می‌گردد.
سپس، الگوی کلّیِ حرکت به‌سوی تمرکز قدرت تحلیل و با یک نمودار راداری
(\en{radar chart}) تصویر‌سازی خواهد شد.
\end{infoBox}

%============================================================
\section{مقدمه: ضرورت تحلیل تطبیقی}
\label{sec:ch13-intro}

تا اینجای کتاب، هر بزنگاه به‌صورت مستقل بررسی شد.
اکنون پرسش اصلی این است:
آیا الگوی واحدی در رفتار خمینی قابل شناسایی است؟%
\footnote{برای اهمیت تحلیل الگومحور در تاریخ سیاسی
ر.ک.\ \en{Theda Skocpol},
\textit{\en{States and Social Revolutions}}, 1979, pp.\,33–40.}
\keyword{تحلیل تطبیقی} (\en{comparative analysis})
ابزار اصلیِ پاسخ‌گویی به این پرسش است:
با کنار هم نهادنِ هفت بزنگاه در یک ماتریس واحد،
هم نقاط اشتراک و هم موارد استثنا روشن می‌شوند.%
\footnote{%
\en{Arend Lijphart},
``Comparative Politics and the Comparative Method'',
\textit{\en{American Political Science Review}}, 65(3), 1971, pp.\,682–693.}

روش کار بدین قرار است:
\begin{enumerate}
\item برای هر بزنگاه، \emph{کنش واقعی} خمینی ثبت می‌شود
      (مستند به منابع دست‌اول).
\item پیش‌بینیِ نظریِ هر تز دربارهٔ آن بزنگاه صورت‌بندی می‌گردد.
\item درجهٔ سازگاری هر تز با کنش واقعی در مقیاس سه‌درجه‌ای
      (ضعیف / متوسط / قوی) ارزیابی می‌شود.
\item در پایان، جمع‌بندیِ بصری با نمودار ارائه خواهد شد.
\end{enumerate}

%============================================================
\section{تعریف معیارهای سازگاری}
\label{sec:ch13-criteria}

جدول~\ref{tab:fit-scale} مقیاس سه‌درجه‌ای به‌کاررفته را نشان می‌دهد.

\begin{table}[ht]
\centering
\caption{مقیاس ارزیابیِ سازگاریِ تز با شواهد}
\label{tab:fit-scale}
\begin{tabular}{>{\centering\arraybackslash}m{2.8cm}
                >{\arraybackslash}m{9.5cm}}
\toprule
\textbf{درجه} & \textbf{تعریف عملیاتی} \\
\midrule
\cellcolor{cGreen!15}\textbf{قوی (۳)}
  & کنش خمینی دقیقاً با پیش‌بینی تز هم‌خوان است؛
    شواهد مستقیم در سخنرانی‌ها یا اسناد موجود است. \\
\addlinespace
\cellcolor{cGold!15}\textbf{متوسط (۲)}
  & تز می‌تواند کنش را توضیح دهد اما نیاز به فروض کمکی دارد؛
    شواهد غیرمستقیم‌اند. \\
\addlinespace
\cellcolor{cAccent!15}\textbf{ضعیف (۱)}
  & تز ناچار به توجیه‌های ثانوی است؛
    شواهد بیشتر ناسازگار با پیش‌بینی‌اند تا سازگار. \\
\bottomrule
\end{tabular}
\end{table}

این مقیاس، هرچند ساده‌سازی‌شده،
امکان مقایسهٔ نظام‌مند را فراهم می‌آورد.%
\footnote{%
دربارهٔ محدودیت‌های مقیاس‌گذاری کیفی
ر.ک.\ \en{Alexander L.\ George \& Andrew Bennett},
\textit{\en{Case Studies and Theory Development in the Social Sciences}},
2005, ch.\,8.}

%============================================================
\section{جدول اصلیِ تطبیقی: هفت بزنگاه × چهار تز}
\label{sec:ch13-master}

جدول~\ref{tab:master-matrix} ماتریس اصلی را ارائه می‌دهد.
ستون نخست بزنگاه و فصل مربوطه را نشان می‌دهد؛
ستون دوم خلاصهٔ کنش واقعی خمینی را ثبت می‌کند؛
و چهار ستون بعدی نمرهٔ سازگاری هر تز را ارائه می‌کنند.
ستون آخر مشخص می‌سازد کدام تز بیشترین سازگاری را دارد.

{%
\small
\begin{longtable}{%
  >{\centering\arraybackslash}m{2.4cm}
  >{\arraybackslash}m{3.2cm}
  >{\centering\arraybackslash}m{1.3cm}
  >{\centering\arraybackslash}m{1.3cm}
  >{\centering\arraybackslash}m{1.3cm}
  >{\centering\arraybackslash}m{1.3cm}
  >{\centering\arraybackslash}m{1.8cm}%
}
\caption{ماتریس تطبیقی: هفت بزنگاه × چهار تز}
\label{tab:master-matrix} \\
\toprule
\textbf{بزنگاه (فصل)}
  & \textbf{کنش واقعی خمینی}
  & \cellcolor{cGreen!20}\textbf{ت‌۱}
  & \cellcolor{cAccent!20}\textbf{ت‌۲}
  & \cellcolor{cGold!20}\textbf{ت‌۳}
  & \cellcolor{cConsolid!20}\textbf{ت‌۴}
  & \textbf{تز برتر} \\
\midrule
\endfirsthead

\caption[]{(ادامه)} \\
\toprule
\textbf{بزنگاه (فصل)}
  & \textbf{کنش واقعی خمینی}
  & \cellcolor{cGreen!20}\textbf{ت‌۱}
  & \cellcolor{cAccent!20}\textbf{ت‌۲}
  & \cellcolor{cGold!20}\textbf{ت‌۳}
  & \cellcolor{cConsolid!20}\textbf{ت‌۴}
  & \textbf{تز برتر} \\
\midrule
\endhead

\midrule
\multicolumn{7}{r}{\footnotesize \textit{ادامه در صفحهٔ بعد}\,\dots} \\
\endfoot

\bottomrule
\endlastfoot

%--- ROW 1 -------------------------------------------------
وعده‌های پاریس (فصل~۶)
  & تصریح بر عدم تصدّی مقام؛
    تأکید بر نقش هدایتی–معنوی%
    \footnote{صحیفهٔ امام، ج\,۵، صص\,۳۸۸–۳۹۰؛ مصاحبه با
    \en{Le Monde}، ۱۶ دی ۱۳۵۷.}
  & \cellcolor{cGreen!15} ۳
  & \cellcolor{cAccent!15} ۳
  & \cellcolor{cGold!15} ۲
  & \cellcolor{cConsolid!15} ۲
  & ت‌۱ / ت‌۲ \\
\addlinespace

%--- ROW 2 -------------------------------------------------
برخورد با بختیار (فصل~۷)
  & نصب بازرگان؛ ایجاد ساختار موازی؛ کنارزدن بختیار%
    \footnote{بازرگان، \textit{انقلاب ایران در دو حرکت}، ص\,۷۲؛
    صحیفهٔ امام، ج\,۶، ص\,۵۴.}
  & \cellcolor{cGreen!15} ۱
  & \cellcolor{cAccent!15} ۳
  & \cellcolor{cGold!15} ۲
  & \cellcolor{cConsolid!15} ۲
  & \cellcolor{cAccent!10} ت‌۲ \\
\addlinespace

%--- ROW 3 -------------------------------------------------
مجلس خبرگان / قانون اساسی (فصل~۸)
  & طرح ولایت فقیه مطلقه؛
    بازنویسی پیش‌نویس%
    \footnote{صورت مشروح مذاکرات مجلس خبرگان قانون اساسی، ۱۳۵۸،
    جلسات ۳۵–۶۷؛
    \en{Schirazi},
    \textit{\en{The Constitution of Iran}}, 1997, pp.\,22–48.}
  & \cellcolor{cGreen!15} ۱
  & \cellcolor{cAccent!15} ۳
  & \cellcolor{cGold!15} ۳
  & \cellcolor{cConsolid!15} ۲
  & \cellcolor{cAccent!10} ت‌۲ / ت‌۳ \\
\addlinespace

%--- ROW 4 -------------------------------------------------
ترورهای فرقان (فصل~۹)
  & بهره‌گیری از بحران امنیتی
    برای تحکیم نهادهای انقلابی%
    \footnote{%
    \en{Abrahamian},
    \textit{\en{Radical Islam}}, 1989, pp.\,186–193.}
  & \cellcolor{cGreen!15} ۲
  & \cellcolor{cAccent!15} ۲
  & \cellcolor{cGold!15} ۲
  & \cellcolor{cConsolid!15} ۳
  & \cellcolor{cConsolid!10} ت‌۴ \\
\addlinespace

%--- ROW 5 -------------------------------------------------
بحران گروگان‌گیری (فصل~۱۰)
  & تأیید ضمنی و سپس صریح؛
    ابزارسازی از بحران برای رفراندوم قانون اساسی و حذف رقبا%
    \footnote{%
    \en{Bowden},
    \textit{\en{Guests of the Ayatollah}}, 2006, pp.\,204–219;
    صحیفهٔ امام، ج\,۱۱، ص\,۱۲۳.}
  & \cellcolor{cGreen!15} ۱
  & \cellcolor{cAccent!15} ۳
  & \cellcolor{cGold!15} ۲
  & \cellcolor{cConsolid!15} ۳
  & \cellcolor{cAccent!10} ت‌۲ / ت‌۴ \\
\addlinespace

%--- ROW 6 -------------------------------------------------
جنگ ایران و عراق (فصل~۱۱)
  & رد پیشنهادهای صلح اولیه؛
    تبدیل جنگ به ابزار بسیج داخلی%
    \footnote{خاطرات هاشمی رفسنجانی، \textit{دوران دفاع مقدس}، ج\,۱،
    صص\,۱۱۲–۱۲۸؛
    \en{Crist},
    \textit{\en{The Twilight War}}, 2012, pp.\,88–103.}
  & \cellcolor{cGreen!15} ۱
  & \cellcolor{cAccent!15} ۲
  & \cellcolor{cGold!15} ۳
  & \cellcolor{cConsolid!15} ۳
  & \cellcolor{cGold!10} ت‌۳ / ت‌۴ \\
\addlinespace

%--- ROW 7 -------------------------------------------------
حذف مجاهدین / بنی‌صدر / تصفیه‌ها / کشتار ۱۳۶۷ (فصل~۱۲)
  & صدور فتوای اعدام‌های جمعی؛
    بالاترین درجهٔ تمرکز قدرت%
    \footnote{خاطرات آیت‌الله منتظری (نسخهٔ اینترنتی)، نامه‌های ۱۳۶۷؛
    \en{Robertson},
    \textit{\en{The Massacre of Political Prisoners in Iran, 1988}},
    2010, pp.\,54–87;
    \en{Abrahamian},
    \textit{\en{Tortured Confessions}}, 1999, pp.\,209–228.}
  & \cellcolor{cGreen!15} ۱
  & \cellcolor{cAccent!15} ۳
  & \cellcolor{cGold!15} ۲
  & \cellcolor{cConsolid!15} ۳
  & \cellcolor{cAccent!10} ت‌۲ / ت‌۴ \\

\end{longtable}
}% end \small

%============================================================
\section{خوانش جدول: الگوهای بارز}
\label{sec:ch13-patterns}

\subsection{الگوی اول: حرکت یک‌سویه به‌سوی تمرکز}

با مرور ستون «کنش واقعی» از بالا به پایین،
مشاهده می‌شود که در هر بزنگاه جدید،
خمینی \keyword{قدرت بیشتری} بر‌ عهده می‌گیرد یا تحکیم می‌کند.%
\footnote{%
\en{Milani},
\textit{\en{The Making of Iran's Islamic Revolution}},
2nd ed., 1994, pp.\,232–234
نیز بر همین الگوی تمرکز تأکید دارد.}
هیچ بزنگاهی وجود ندارد که در آن خمینی قدرتی را واگذار کرده باشد
یا نهادی مستقل را تقویت کرده باشد بدون آنکه آن نهاد
زیر نظارت مستقیم ولایت فقیه درآید.

این الگو برای
\textcolor{cGreen}{تز اول (صداقت + اضطرار)}
مشکل‌آفرین است:
اگر قدرت‌گیری هر بار فقط واکنشی به بحران بود،
باید دست‌کم در یکی از بزنگاه‌ها پس از رفع بحران
بازگشتی به موقعیت قبلی می‌دیدیم.%
\footnote{%
بشیریه، \textit{دیباچه‌ای بر جامعه‌شناسی سیاسی ایران}، ص\,۱۸۵
به همین نکته اشاره می‌کند.}
چنین بازگشتی هرگز رخ نداد.

\subsection{الگوی دوم: بهره‌گیری نظام‌مند از بحران}

در حداقل چهار بزنگاه
(بختیار، گروگان‌گیری، ترورهای فرقان، جنگ)
خمینی از بحران بیرونی یا امنیتی
برای پیشبرد دستور کار داخلی بهره برد.%
\footnote{گنجی، \textit{تلقی فاشیستی از دین و حکومت}، صص\,۱۲۰–۱۳۵.}
این الگو هم با \textcolor{cAccent}{تز دوم (فریب آگاهانه)} و
هم با \textcolor{cConsolid}{تز چهارم (فشار نخبگان)} سازگار است،
اما با تز اول (که بحران‌ها را \emph{علتِ} تغییر مسیر می‌داند)
تنها به‌طور ناقص انطباق دارد.

\subsection{الگوی سوم: نقطهٔ عطف قانون اساسی}

بزنگاه مجلس خبرگان
(ردیف ۳ جدول~\ref{tab:master-matrix})
جایی است که هم \textcolor{cAccent}{تز دوم} و هم
\textcolor{cGold}{تز سوم (تطور ایدئولوژیک)}
بالاترین نمره (۳) را می‌گیرند.
این همان لحظه‌ای است که وعدهٔ نوفل‌لوشاتو
به‌صورت نهادی و حقوقی نقض می‌شود.%
\footnote{%
\en{Arjomand},
\textit{\en{The Turban for the Crown}}, 1988, pp.\,135–148.}

\subsection{الگوی چهارم: ضعف مستمر تز اول}

تز اول در پنج بزنگاه از هفت بزنگاه
پایین‌ترین نمره (۱) را دریافت می‌کند.
تنها در بزنگاه اول (وعده‌های پاریس)
نمرهٔ ۳ و در بزنگاه ترورهای فرقان نمرهٔ ۲ کسب می‌کند.
بنابراین، تز اول بیشتر ناظر بر نیت اولیه است
تا بر رفتار واقعی پس از پیروزی انقلاب.%
\footnote{کدیور، \textit{حکومت ولایی}، ص\,۲۱۴
تمایز میان نیت اولیه و عملکرد نهایی را محوری می‌داند.}

%============================================================
\section{جمع‌بندی عددی: مجموع نمرات هر تز}
\label{sec:ch13-scores}

جدول~\ref{tab:score-summary} جمع نمرات هر تز در هفت بزنگاه
و میانگین آن‌ها را نشان می‌دهد.

\begin{table}[ht]
\centering
\caption{جمع و میانگین نمرات سازگاری هر تز (از مجموع ۲۱)}
\label{tab:score-summary}
\begin{tabular}{>{\centering\arraybackslash}m{4.2cm}
                >{\centering\arraybackslash}m{2.2cm}
                >{\centering\arraybackslash}m{2.2cm}}
\toprule
\textbf{تز} & \textbf{مجموع نمره (از ۲۱)} & \textbf{میانگین} \\
\midrule
\cellcolor{cGreen!15}
  تز ۱: صداقت + اضطرار & ۱۰ & ۱٫۴۳ \\
\cellcolor{cAccent!15}
  تز ۲: فریب آگاهانه & ۱۹ & ۲٫۷۱ \\
\cellcolor{cGold!15}
  تز ۳: تطور ایدئولوژیک & ۱۶ & ۲٫۲۹ \\
\cellcolor{cConsolid!15}
  تز ۴: فشار نخبگان & ۱۸ & ۲٫۵۷ \\
\bottomrule
\end{tabular}
\end{table}

بدین ترتیب ترتیب سازگاری کلی چنین است:
\[
\textcolor{cAccent}{\text{ت‌۲}}\;(۱۹)
\;>\;
\textcolor{cConsolid}{\text{ت‌۴}}\;(۱۸)
\;>\;
\textcolor{cGold}{\text{ت‌۳}}\;(۱۶)
\;>\;
\textcolor{cGreen}{\text{ت‌۱}}\;(۱۰)
\]

%============================================================
\section{نمودار راداری: پروفایل هر تز در هفت بزنگاه}
\label{sec:ch13-radar}

نمودار زیر نمرهٔ هر تز را در هر بزنگاه
بر روی یک محور جداگانه ترسیم می‌کند.%
\footnote{%
نمودارهای راداری برای مقایسهٔ چند‌بُعدیِ تزها
ابزار رایج در علوم اجتماعیِ کمّی‌اند.
ر.ک.\ \en{Edward Tufte},
\textit{\en{The Visual Display of Quantitative Information}},
2nd ed., 2001, p.\,178.}

\begin{figure}[ht]
\centering
\begin{tikzpicture}[scale=1.1]
  % --- helpers ------------------------------------------
  \def\nAxes{7}
  \def\maxVal{3}
  \pgfmathsetmacro{\angleStep}{360/\nAxes}

  % axis labels (Persian, shortened)
  \def\axisLabels{%
    {پاریس},
    {بختیار},
    {خبرگان},
    {فرقان},
    {گروگان‌گیری},
    {جنگ},
    {تصفیه/۶۷}%
  }

  % data: T1, T2, T3, T4 (in order of 7 turning points)
  \def\dataA{3,1,1,2,1,1,1}     % T1  (green)
  \def\dataB{3,3,3,2,3,2,3}     % T2  (red)
  \def\dataC{2,2,3,2,2,3,2}     % T3  (gold)
  \def\dataD{2,2,2,3,3,3,3}     % T4  (purple)

  % --- background grid ----------------------------------
  \foreach \r in {1,2,3} {
    \draw[cGray!30, thin]
      \foreach \i [count=\idx from 0] in {1,...,\nAxes} {%
        ({90-\idx*\angleStep}:{\r*1.1}) --
      } cycle;
  }

  % --- axes ---------------------------------------------
  \foreach \i [count=\idx from 0] in {1,...,\nAxes} {
    \draw[cGray!50] (0,0) -- ({90-\idx*\angleStep}:{\maxVal*1.1+0.4});
  }

  % --- axis labels --------------------------------------
  \foreach \lbl [count=\idx from 0] in \axisLabels {
    \pgfmathsetmacro{\ang}{90-\idx*\angleStep}
    \node[font=\scriptsize, align=center]
      at ({\ang}:{\maxVal*1.1+0.95}) {\lbl};
  }

  % --- macro: draw one data polygon ---------------------
  \newcommand{\drawPoly}[3]{% #1=data list, #2=draw color, #3=fill color
    \foreach \val [count=\idx from 0] in #1 {
      \pgfmathsetmacro{\ang}{90-\idx*\angleStep}
      \pgfmathsetmacro{\radius}{\val*1.1}
      \coordinate (P\idx) at ({\ang}:{\radius});
    }
    \draw[#2, thick, fill=#3, fill opacity=0.12]
      (P0) \foreach \idx in {1,...,6}{ -- (P\idx) } -- cycle;
  }

  % --- plot polygons ------------------------------------
  \drawPoly{\dataA}{cGreen}{cGreen}
  \drawPoly{\dataB}{cAccent}{cAccent}
  \drawPoly{\dataC}{cGold}{cGold}
  \drawPoly{\dataD}{cConsolid}{cConsolid}

  % --- legend -------------------------------------------
  \node[anchor=north west, font=\scriptsize] at (4.2,3.6) {%
    \begin{tabular}{cl}
      \tikz\draw[cGreen, thick, fill=cGreen!15] (0,0) rectangle (0.35,0.2); & تز ۱ \\
      \tikz\draw[cAccent, thick, fill=cAccent!15] (0,0) rectangle (0.35,0.2); & تز ۲ \\
      \tikz\draw[cGold, thick, fill=cGold!15] (0,0) rectangle (0.35,0.2); & تز ۳ \\
      \tikz\draw[cConsolid, thick, fill=cConsolid!15] (0,0) rectangle (0.35,0.2); & تز ۴ \\
    \end{tabular}
  };
\end{tikzpicture}
\caption{نمودار راداری: سازگاری هر تز با هفت بزنگاه}
\label{fig:radar}
\end{figure}

چنان‌که نمودار~\ref{fig:radar} نشان می‌دهد،
\textcolor{cAccent}{تز دوم} و \textcolor{cConsolid}{تز چهارم}
بیشترین پوشش سطحی را دارند،
در حالی که \textcolor{cGreen}{تز اول} کوچک‌ترین چندضلعی
را تشکیل می‌دهد و تنها در محور «پاریس» به حداکثر می‌رسد.

%============================================================
\section{پرسش کلیدی: آیا مسیر قدرت همیشه یک‌سویه بود؟}
\label{sec:ch13-direction}

برای پاسخ دقیق‌تر، شاخص ساده‌ای تعریف می‌کنیم:

\begin{noteBox}[شاخص جهت قدرت]
اگر $\Delta_i$ نشان‌دهندهٔ تغییر میزان قدرت نهادیِ خمینی
در بزنگاه $i$ نسبت به بزنگاه $i{-}1$ باشد،
جهت‌گیری کلی را با علامتِ $\sum_{i=1}^{7}\Delta_i$ نشان می‌دهیم.
مقدار مثبت = تمرکز؛ مقدار منفی = واگذاری.
\end{noteBox}

در هیچ‌یک از بزنگاه‌ها $\Delta_i < 0$ نیست.
بنابراین، $\sum \Delta_i > 0$ و مسیر قدرت
\keyword{یک‌سویه و انباشتی} بوده است.%
\footnote{%
\en{Bakhash},
\textit{\en{The Reign of the Ayatollahs}}, 1984, pp.\,52–73
همین انباشت تدریجی قدرت را
مهم‌ترین ویژگی دوران ۱۳۵۸–۱۳۶۷ می‌داند.}

%============================================================
\section{تکمیل تحلیل: نمودار روند تمرکز قدرت}
\label{sec:ch13-trend}

\begin{figure}[ht]
\centering
\begin{tikzpicture}
  \begin{axis}[
      width=13cm, height=6.5cm,
      xlabel={\scriptsize بزنگاه},
      ylabel={\scriptsize شاخص تجمیعی قدرت},
      xtick={1,...,7},
      xticklabels={پاریس, بختیار, خبرگان, فرقان, گروگان, جنگ, تصفیه/۶۷},
      xticklabel style={rotate=35, anchor=east, font=\tiny},
      ymin=0, ymax=22,
      grid=major,
      grid style={cGray!20},
      legend style={at={(0.02,0.98)}, anchor=north west, font=\tiny},
      thick,
    ]
    % cumulative "power index" (illustrative)
    \addplot[cAccent, mark=square*, very thick]
      coordinates {(1,0)(2,4)(3,9)(4,12)(5,16)(6,18)(7,21)};
    \addlegendentry{شاخص تمرکز قدرت}

    % hypothetical "promise compliance" declining
    \addplot[cGreen, mark=triangle*, dashed, thick]
      coordinates {(1,20)(2,15)(3,8)(4,7)(5,4)(6,3)(7,0)};
    \addlegendentry{شاخص پای‌بندی به وعده}
  \end{axis}
\end{tikzpicture}
\caption{روند متقابل تمرکز قدرت و پای‌بندی به وعده (شِماتیک)}
\label{fig:trend-power}
\end{figure}

نمودار~\ref{fig:trend-power} به‌صورت شماتیک نشان می‌دهد
که دو خط — شاخص تمرکز قدرت و شاخص پای‌بندی به وعده —
تقریباً \keyword{آینه‌ای} (\en{mirror-image})
رفتار می‌کنند:
هر افزایش در اولی با کاهش متناظری در دومی همراه است.%
\footnote{این نمودار تنها جنبهٔ تصویری دارد
و بر اساس ارزیابی کیفی نگارنده ترسیم شده است.
هدف از آن تسهیل درک بصری الگو است، نه ارائهٔ دادهٔ کمّی دقیق.}

%============================================================

\begin{principleBox}[خلاصهٔ فصل ۱۳]
\begin{itemize}
\item جدول تطبیقیِ هفت بزنگاه × چهار تز نشان داد که
      \textcolor{cAccent}{تز دوم (فریب آگاهانه)} با مجموع نمرهٔ ۱۹ از ۲۱
      بیشترین سازگاری را با شواهد دارد.
\item \textcolor{cConsolid}{تز چهارم (فشار نخبگان)} با ۱۸
      و \textcolor{cGold}{تز سوم (تطور ایدئولوژیک)} با ۱۶
      در رتبه‌های بعدی قرار دارند.
\item \textcolor{cGreen}{تز اول (صداقت + اضطرار)} با ۱۰
      ضعیف‌ترین عملکرد تطبیقی را دارد
      و بیشتر ناظر بر بزنگاه نخست (پاریس) است.
\item الگوی حرکت قدرت یک‌سویه و انباشتی است؛
      هیچ بزنگاه بازگشتی مشاهده نشد.
\item نمودار راداری و نمودار روند تمرکز
      هر دو مؤید همین نتیجه‌اند.
\end{itemize}
\end{principleBox}


%%%%%%%%%%%%%%%%%%%%%%%%%%%%%%%%%%%%%%%%%%%%%%%%%%%%%%%%%%%%
%%  CHAPTER 14
%%%%%%%%%%%%%%%%%%%%%%%%%%%%%%%%%%%%%%%%%%%%%%%%%%%%%%%%%%%%
\chapter{خوانش‌های رقیب در بوتهٔ نقد}
\label{ch:mutual-critique}

%------------------------------------------------------------
\begin{infoBox}[چکیدهٔ فصل ۱۴]
در این فصل، هر یک از شش جفت ممکن از چهار تز
(\en{$\binom{4}{2}=6$ pairs})
در برابر یکدیگر نقد می‌شوند.
نخست نقاط قوت و ضعف هر تز فهرست‌وار ارائه می‌شود،
سپس نقد متقابل هر جفت بررسی می‌گردد.
در پایان، \keyword{ماتریس سازگاری} نشان خواهد داد
کدام تزها با یکدیگر قابل ترکیب‌اند و کدام‌ها متناقض.
\end{infoBox}

%============================================================
\section{نقاط قوت و ضعف هر تز: نگاه اجمالی}
\label{sec:ch14-strengths}

جدول~\ref{tab:strength-weakness} خلاصه‌ای از مهم‌ترین
نقاط قوت و ضعف هر تز را ارائه می‌دهد.%
\footnote{این ارزیابی بر پایهٔ تحلیل فصل‌های ۳–۱۳ است
و معیارها عبارت‌اند از:
(الف) سازگاری با شواهد اولیه،
(ب) قدرت پیش‌بینی‌کنندگی،
(ج) پارسیمونی (\en{parsimony}).}

{%
\small
\begin{longtable}{%
  >{\centering\arraybackslash}m{2.5cm}
  >{\arraybackslash}m{5cm}
  >{\arraybackslash}m{5cm}}
\caption{نقاط قوت و ضعف هر تز}
\label{tab:strength-weakness} \\
\toprule
\textbf{تز} & \textbf{نقاط قوت} & \textbf{نقاط ضعف} \\
\midrule
\endfirsthead
\caption[]{(ادامه)} \\
\toprule
\textbf{تز} & \textbf{نقاط قوت} & \textbf{نقاط ضعف} \\
\midrule
\endhead
\bottomrule
\endlastfoot

\cellcolor{cGreen!12}
\textbf{ت‌۱:\ صداقت\newline + اضطرار}
  & \begin{itemize}[nosep, leftmargin=*]
    \item با لحن مصاحبه‌های ۱۳۵۷ سازگار%
      \footnote{صحیفهٔ امام، ج\,۵، صص\,۳۸۸–۳۹۰.}
    \item به پیچیدگیِ روان‌شناختی رهبران اذعان دارد
    \item شواهد برخی تردیدهای خمینی در خاطرات یاران%
      \footnote{یزدی، \textit{آخرین تلاش‌ها در آخرین روزها}، ص\,۱۸۰.}
    \end{itemize}
  & \begin{itemize}[nosep, leftmargin=*]
    \item در ۵ از ۷ بزنگاه ضعیف‌ترین سازگاری (فصل~۱۳)
    \item عدم بازگشت قدرت پس از رفع بحران
    \item نمی‌تواند فتوای ۱۳۶۷ را توجیه کند%
      \footnote{\en{Robertson}, \textit{\en{Massacre}}, 2010, p.\,122.}
    \end{itemize} \\
\addlinespace

\cellcolor{cAccent!12}
\textbf{ت‌۲:\ فریب\newline آگاهانه}
  & \begin{itemize}[nosep, leftmargin=*]
    \item بالاترین نمرهٔ تطبیقی (۱۹ از ۲۱)
    \item با آموزهٔ تقیه قابل ارتباط%
      \footnote{\en{Moin}, \textit{\en{Khomeini}}, 1999, pp.\,200–204.}
    \item شواهد مستقیم از نامه‌های محرمانه%
      \footnote{خاطرات منتظری، بخش مکاتبات ۱۳۶۱–۱۳۶۴.}
    \end{itemize}
  & \begin{itemize}[nosep, leftmargin=*]
    \item فرض ذهن \keyword{یکپارچهٔ} (\en{unitary mind}) محل تردید است%
      \footnote{\en{Jon Elster},
      \textit{\en{Sour Grapes}}, 1983, pp.\,141–166.}
    \item نمی‌تواند آثار فقهی پیش‌ از ۱۳۵۷ را نادیده بگیرد%
      \footnote{\textit{ولایت فقیه} (نجف، ۱۳۴۸ ش.)
      که ایدهٔ حکومت اسلامی در آن مطرح شده بود.}
    \item خطر افتادن در \en{conspiracy theory}
    \end{itemize} \\
\addlinespace

\cellcolor{cGold!12}
\textbf{ت‌۳:\ تطور\newline ایدئولوژیک}
  & \begin{itemize}[nosep, leftmargin=*]
    \item تغییر در آثار فقهی خمینی مستند است%
      \footnote{کدیور، \textit{نظریه‌های دولت در فقه شیعه}، صص\,۱۵۲–۱۸۰.}
    \item از نظر روان‌شناختی واقع‌بینانه‌تر از ت‌۲
    \item مفهوم مصلحت نظام سیر تطوری دارد%
      \footnote{صحیفهٔ امام، ج\,۲۱، نامهٔ معروف ۱۳۶۶
      دربارهٔ تقدم مصلحت نظام بر احکام اولیه.}
    \end{itemize}
  & \begin{itemize}[nosep, leftmargin=*]
    \item نمی‌تواند سرعت تمرکز ۱۳۵۸ را توضیح دهد
    \item کتاب \textit{ولایت فقیه} (۱۳۴۸) نشان‌دهندهٔ
      وجود ایدهٔ حاکمیت فقیه \emph{پیش‌ از} انقلاب
    \item مرز میان «تطور» و «فرصت‌طلبی» مبهم
    \end{itemize} \\
\addlinespace

\cellcolor{cConsolid!12}
\textbf{ت‌۴:\ فشار\newline نخبگان}
  & \begin{itemize}[nosep, leftmargin=*]
    \item نقش بهشتی، رفسنجانی، خامنه‌ای مستند%
      \footnote{خاطرات رفسنجانی، ج\,۱، صص\,۸۵–۱۱۰.}
    \item با تحلیل‌های جامعه‌شناسی نخبگان سازگار%
      \footnote{\en{Mosca},
      \textit{\en{The Ruling Class}}, 1939;
      \en{Pareto}, \textit{\en{Mind and Society}}, 1935.}
    \item توزیع مسئولیت واقع‌بینانه‌تر از تمرکز بر یک فرد
    \end{itemize}
  & \begin{itemize}[nosep, leftmargin=*]
    \item کاهش مسئولیت فردی خمینی ناسازگار
      با قدرت تصمیم‌گیری‌اش
    \item فتوای ۱۳۶۷ قابل تقلیل به فشار گروهی نیست%
      \footnote{منتظری، خاطرات، نامهٔ ۹ مرداد ۱۳۶۷
      که خمینی شخصاً از او ناراضی شد.}
    \item اسناد «فشار» بیشتر غیرمستقیم‌اند
    \end{itemize} \\

\end{longtable}
}% end \small

%============================================================
\section{نقد متقابل شش جفت تز}
\label{sec:ch14-pairs}

\subsection{جفت ۱: تز ۱ ↔ تز ۲ — صداقت در برابر فریب}
\label{subsec:pair-t1t2}

این بنیادی‌ترین تقابل است:
آیا خمینی در پاریس \keyword{راست} می‌گفت یا \keyword{دروغ}؟

\begin{warningBox}[نقد تز ۱ از منظر تز ۲]
اگر خمینی واقعاً صادق بود،
چرا هیچ‌گاه پس از رفع هر بحران
قدرت را محدود نکرد؟%
\footnote{گنجی، \textit{تلقی فاشیستی}، ص\,۱۴۱.}
الگوی «بحران → تمرکز → عدم بازگشت»
بیش از آنکه با اضطرار سازگار باشد
با طرح از‌پیش‌اندیشیده هم‌خوان است.
\end{warningBox}

\begin{warningBox}[نقد تز ۲ از منظر تز ۱]
تز فریب آگاهانه فرض می‌کند
خمینی در ۱۳۵۷ تصویر دقیقی از آیندهٔ ده‌سالهٔ
خود داشته است — فرضی بسیار قوی
که مستلزم \en{perfect foresight} است.%
\footnote{\en{Kurzman},
\textit{\en{The Unthinkable Revolution}}, 2004, pp.\,5–7
بر غیرقابل‌پیش‌بینی بودنِ حتی خودِ انقلاب تأکید دارد.}
مصاحبه‌های متعدد پاریس لحنی دارند
که جعل آن در آن مقیاس بسیار دشوار بوده است.
\end{warningBox}

\subsection{جفت ۲: تز ۱ ↔ تز ۳ — صداقت در برابر تطور}
\label{subsec:pair-t1t3}

\begin{noteBox}[همپوشانی ممکن]
تز اول و سوم در یک نقطه مشترک‌اند:
هر دو خمینی را فاعل اخلاقی‌ای می‌بینند که
نیتش در ۱۳۵۷ واقعاً عدم حاکمیت بوده است.
تفاوت در تبیین علت تغییر است:
ت‌۱ می‌گوید بحران بیرونی؛
ت‌۳ می‌گوید تحول درونیِ فکری.%
\footnote{سروش، \textit{فربه‌تر از ایدئولوژی}، ص\,۹۵
دربارهٔ تغییر پارادایم درون‌ایدئولوژیک.}
\end{noteBox}

نقد تز ۳ بر تز ۱:
«اضطرار» بسیار ساده‌انگارانه است؛
تحول اندیشه فرایندی تدریجی و مستند است،
نه صرفاً واکنش لحظه‌ای.%
\footnote{کدیور، \textit{حکومت ولایی}، ص\,۱۹۵.}

نقد تز ۱ بر تز ۳:
اگر تطور واقعاً وجود داشت،
سرعت عمل ۱۳۵۸ (فقط چند ماه بعد از وعده)
با فرایند تدریجیِ تطور ایدئولوژیک ناسازگار است.%
\footnote{\en{Milani},
\textit{\en{Making of Iran's Islamic Revolution}}, 1994, p.\,196.}

\subsection{جفت ۳: تز ۱ ↔ تز ۴ — صداقت در برابر فشار نخبگان}
\label{subsec:pair-t1t4}

نقد تز ۴ بر تز ۱:
اضطرار یک نفره نبود؛ ساختار نخبگان انقلابی
اصلاً اجازهٔ واگذاری قدرت را نمی‌داد.%
\footnote{حجاریان، \textit{از شاهد قدسی تا شاهد بازاری}، ص\,۲۰۸.}

نقد تز ۱ بر تز ۴:
اگر فشار نخبگان تعیین‌کننده بود،
چرا خمینی توانست تمام نخبگان مخالف
(بازرگان، بنی‌صدر، منتظری) را کنار بزند؟%
\footnote{بنی‌صدر، \textit{خیانت به امید}، صص\,۳۰۲–۳۱۵.}
این خود نشان‌دهندهٔ ارادهٔ مستقل و قدرتمند رهبر است.

\subsection{جفت ۴: تز ۲ ↔ تز ۳ — فریب در برابر تطور}
\label{subsec:pair-t2t3}

این جفت از نظر منطقی مهم‌ترین تقابل است،
زیرا مشخص می‌کند آیا خمینی
\emph{از ابتدا} قصد حاکمیت داشت (ت‌۲)
یا در مسیر به آن رسید (ت‌۳).

\begin{warningBox}[نقد تز ۲ بر تز ۳]
کتاب \textit{ولایت فقیه} (نجف، ۱۳۴۸ ش.)
به‌روشنی ایدهٔ حکومت مستقیم فقیه را مطرح می‌کند —
یعنی ایده \emph{قبل} از وعده‌های پاریس وجود داشته
و «تطور» در دههٔ ۱۳۵۰ اتفاق افتاده نه ۱۳۵۸.%
\footnote{%
\en{Martin},
\textit{\en{Creating an Islamic State}}, 2000, pp.\,78–95.}
بنابراین، وعده‌های پاریس نمی‌توانند
بازتاب‌دهندهٔ اندیشهٔ واقعی خمینی باشند.
\end{warningBox}

\begin{warningBox}[نقد تز ۳ بر تز ۲]
فریب آگاهانه مستلزم نظریهٔ ذهن
(\en{Theory of Mind}) بسیار قوی
و کنترل کامل بر محیط ارتباطی است.
خمینی در پاریس تحت فشار ده‌ها خبرنگار بود
و وعده‌ها فقط شفاهی نبودند بلکه مکتوب نیز شدند.%
\footnote{صحیفهٔ امام، ج\,۴، پیام «رأی من رأی ملت است»، بهمن ۱۳۵۷.}
واقع‌بینانه‌تر آن است که خمینی
در آن مقطع واقعاً باور داشت اما بعداً تغییر کرد.
\end{warningBox}

\subsection{جفت ۵: تز ۲ ↔ تز ۴ — فریب در برابر فشار نخبگان}
\label{subsec:pair-t2t4}

\begin{noteBox}[سازگاری نسبی]
این دو تز بیش از سایر جفت‌ها قابل ترکیب‌اند:
ممکن است خمینی از ابتدا قصد قدرت داشته (ت‌۲)
\emph{و} نخبگان نیز مستقلاً این مسیر را تقویت کرده باشند (ت‌۴).%
\footnote{\en{Abrahamian},
\textit{\en{Khomeinism}}, 1993, pp.\,17–38.}
\end{noteBox}

نقد تز ۴ بر تز ۲:
فریب کامل نیازی به نخبگان ندارد — اما واقعیت این است
که بدون بهشتی و شورای انقلاب،
تمرکز قدرت در آن سرعت ممکن نبود.%
\footnote{خاطرات رفسنجانی، ج\,۱، صص\,۱۴۸–۱۵۵.}

نقد تز ۲ بر تز ۴:
اگر نخبگان عامل اصلی بودند،
چرا خمینی توانست هر یک را جداگانه حذف کند
بدون آنکه سیستم فرو بپاشد؟
رفسنجانی خود اعتراف کرده
«کلمهٔ آخر را همیشه امام می‌گفتند.»%
\footnote{خاطرات رفسنجانی، ج\,۳، ص\,۲۷۱.}

\subsection{جفت ۶: تز ۳ ↔ تز ۴ — تطور در برابر فشار نخبگان}
\label{subsec:pair-t3t4}

نقد تز ۴ بر تز ۳:
تطور ایدئولوژیک در خلأ رخ نمی‌دهد —
محیط نخبگان اطراف خمینی
همان عاملی است که تطور را جهت‌دهی کرده است.%
\footnote{بشیریه، \textit{دیباچه‌ای بر جامعه‌شناسی سیاسی ایران}،
ص\,۲۱۲.}

نقد تز ۳ بر تز ۴:
تز فشار نخبگان، خمینی را فاقد ارادهٔ مستقل
فکری ترسیم می‌کند — تصویری که با شخصیت
سرسخت و مصممِ او ناسازگار است.%
\footnote{\en{Moin}, \textit{\en{Khomeini}}, 1999, pp.\,8–11.}

%============================================================
\section{ماتریس سازگاری: کدام تزها قابل ترکیب‌اند؟}
\label{sec:ch14-compat}

جدول~\ref{tab:compat-matrix} سازگاری منطقیِ هر جفت تز
را در سه درجه ارزیابی می‌کند:
\textbf{ناسازگار} (×)، \textbf{سازگاری محدود} (△)،
\textbf{سازگار / قابل ترکیب} (○).%
\footnote{معیار سازگاری: آیا دو تز را می‌توان
هم‌زمان بدون تناقض منطقی پذیرفت؟}

\begin{table}[ht]
\centering
\caption{ماتریس سازگاری شش جفت تز}
\label{tab:compat-matrix}
\begin{tabular}{>{\centering\arraybackslash}m{2.5cm}
                >{\centering\arraybackslash}m{2cm}
                >{\centering\arraybackslash}m{2cm}
                >{\centering\arraybackslash}m{2cm}
                >{\centering\arraybackslash}m{2cm}}
\toprule
 & \cellcolor{cGreen!20}\textbf{ت‌۱}
 & \cellcolor{cAccent!20}\textbf{ت‌۲}
 & \cellcolor{cGold!20}\textbf{ت‌۳}
 & \cellcolor{cConsolid!20}\textbf{ت‌۴} \\
\midrule
\cellcolor{cGreen!20}\textbf{ت‌۱}
  & —
  & \cellcolor{cAccent!10} ×
  & \cellcolor{cGold!10} △
  & \cellcolor{cConsolid!10} △ \\
\cellcolor{cAccent!20}\textbf{ت‌۲}
  & \cellcolor{cAccent!10} ×
  & —
  & \cellcolor{cGold!10} ×
  & \cellcolor{cConsolid!10} ○ \\
\cellcolor{cGold!20}\textbf{ت‌۳}
  & \cellcolor{cGold!10} △
  & \cellcolor{cAccent!10} ×
  & —
  & \cellcolor{cConsolid!10} ○ \\
\cellcolor{cConsolid!20}\textbf{ت‌۴}
  & \cellcolor{cConsolid!10} △
  & \cellcolor{cAccent!10} ○
  & \cellcolor